\documentclass[10pt,oneside]{book}

% Cheat-sheet compile toggle.
% Uncomment the next line for compact cheat-sheet styling.
% \notescheatsheettrue
\newif\ifnotescheatsheet
\notescheatsheetfalse
\providecommand{\notescheatsheettoggle}{0}
\ifnum\notescheatsheettoggle=1\relax
  \notescheatsheettrue
\fi

% Page geometry: narrow left margin, generous right margin for notes.
% Cheatsheet mode uses tight uniform margins with no margin column.
\ifnotescheatsheet
\usepackage[
    paper=letterpaper,
    includeheadfoot,
    left=0.5in,
    right=0.5in,
    top=0.5in,
    bottom=0.5in,
    marginparsep=0pt,
    marginparwidth=0pt,
    headheight=15pt,
    headsep=0.12in,
    footskip=0.28in
]{geometry}
\else
\usepackage[
    paper=letterpaper,
    includeheadfoot,
    left=0.50in,
    right=2.13in,
    top=0.70in,
    bottom=0.85in,
    marginparsep=0.18in,
    marginparwidth=1.66in,
    headheight=15pt,
    headsep=0.18in,
    footskip=0.38in
]{geometry}
\fi

% Typography and microtypography.
\usepackage[T1]{fontenc}
\usepackage[utf8]{inputenc}
%\usepackage[semibold]{libertine}
% Core math and theorem support.
\usepackage{amsmath, amssymb, amsthm, mathtools}
\usepackage{centernot}
%\DeclareMathSizes{10}{9.5}{7}{5}
%\DeclareMathSizes{10.5}{10}{7.25}{5.2}
%\DeclareMathSizes{11}{10.4}{7.6}{5.6}
%\DeclareMathSizes{12}{11.3}{8.1}{6}
\setlength{\jot}{0.75pt}
\setlength{\abovedisplayskip}{3pt plus 1pt minus 1pt}
\setlength{\belowdisplayskip}{3pt plus 1pt minus 1pt}
\setlength{\abovedisplayshortskip}{1.5pt plus 1pt minus 0.5pt}
\setlength{\belowdisplayshortskip}{1.5pt plus 1pt minus 0.5pt}
\usepackage{microtype}
\usepackage{textcomp}
\usepackage{setspace}
\ifnotescheatsheet
\setstretch{1.0}
\else
\setstretch{1.0}
\fi
\makeatletter
\ifnotescheatsheet
\else
\renewcommand\normalsize{%
   \@setfontsize\normalsize{10.5}{12.6}%
   \abovedisplayskip 10\p@ \@plus2\p@ \@minus5\p@
   \abovedisplayshortskip \z@ \@plus3\p@
   \belowdisplayshortskip 6\p@ \@plus3\p@ \@minus3\p@
   \belowdisplayskip \abovedisplayskip
   \let\@listi\@listI}
\normalsize
\ifx\MakeRobust\@undefined \else
    \MakeRobust\normalsize
\fi
\fi
\makeatother
\raggedbottom

\usepackage{dsfont}
\numberwithin{equation}{chapter}
\usepackage{mathrsfs}
\DeclareFontFamily{U}{rsfs}{\skewchar\font127 }
\DeclareFontShape{U}{rsfs}{m}{n}{<-6> rsfs5 <6-8> rsfs7 <8-> rsfs10}{}

% Tables, figures, colors, and graphics.
\usepackage{float}
\usepackage{graphicx}
\usepackage{booktabs}
\usepackage{threeparttable}
\usepackage{array}
\usepackage{tabularx}
\usepackage{tabularray}
\usepackage{caption}
\usepackage{siunitx}
\usepackage{codehigh}
\usepackage[normalem]{ulem}
\usepackage{cancel}
\usepackage{fontawesome5}
\usepackage{xcolor}
\usepackage{tikz}
\usetikzlibrary{calc}
\usepackage{pgfplots}
\pgfplotsset{compat=1.18}
\usepackage{soul}
\UseTblrLibrary{booktabs}
\UseTblrLibrary{siunitx}
\providecommand{\tinytableTabularrayUnderline}[1]{\underline{#1}}
\providecommand{\tinytableTabularrayStrikeout}[1]{\sout{#1}}
\NewTableCommand{\tinytableDefineColor}[3]{\definecolor{#1}{#2}{#3}}

% Boxes and marginalia.
\usepackage[most]{tcolorbox}
\tcbuselibrary{theorems}
\newcommand{\lecboxednoteribbon}{black!22}
\newcommand{\lecsetboxednoteribbon}[1]{\gdef\lecboxednoteribbon{#1}}
\newcommand{\lecresetboxednoteribbon}{\lecsetboxednoteribbon{black!22}}
\ifnotescheatsheet
\tcbset{
  lecrcodeboxstyle/.style={
    colback=white,
    colframe=black,
    boxrule=0.45pt,
    arc=0pt,
    left=4pt,
    right=4pt,
    top=4pt,
    bottom=4pt
  },
  lecnoteboxstyle/.style={
    arc=0pt,
    colback=white,
    colframe=black,
    boxrule=0.25pt,
    left=3pt,
    right=3pt,
    top=3pt,
    bottom=3pt
  },
  lecthmbasetitlestyle/.style={
    arc=0pt,
    left=4pt,
    right=4pt,
    top=3pt,
    bottom=3pt
  },
  lecthmbasestyle/.style={
    arc=0pt,
    left=5pt,
    right=5pt,
    top=5pt,
    bottom=5pt
  }
}
\newcommand{\lecnotefont}{\tiny}
\newcommand{\lecboxednotefootnotefont}{\fontsize{4.5}{6}\selectfont}
\else
\tcbset{
  lecrcodeboxstyle/.style={
    colback=gray!5,
    colframe=black!65,
    boxrule=0.45pt,
    arc=1.4mm,
    left=6pt,
    right=6pt,
    top=6pt,
    bottom=6pt
  },
  lecnoteboxstyle/.style={
    arc=0.8mm,
    outer arc=0.8mm,
    sharp corners=west,
    colback=white,
    colframe=black!45,
    boxrule=0.25pt,
    left=4pt,
    right=4pt,
    top=3pt,
    bottom=3pt
  },
  lecthmbasetitlestyle/.style={
    arc=1.4mm,
    left=6pt,
    right=6pt,
    top=4pt,
    bottom=4pt
  },
  lecthmbasestyle/.style={
    arc=1.4mm,
    left=8pt,
    right=8pt,
    top=8pt,
    bottom=8pt
  }
}
\newcommand{\lecnotefont}{\scriptsize}
\newcommand{\lecboxednotefootnotefont}{\tiny}
\fi
\usepackage{varwidth}
\usepackage{marginnote}
\usepackage{marginfix}
\usepackage{zref-savepos}
\setlength{\marginparpush}{9pt}
\renewcommand*{\marginfont}{\footnotesize}

% Links.
\usepackage{hyperref}
\hypersetup{
    colorlinks=true,
    linkcolor=black,
    citecolor=blue!70!black,
    urlcolor=blue!70!black,
    pdfauthor={Marcus Kuo},
    pdftitle={Notes}
}

% Lists and TOC helpers.
\usepackage{enumitem}
\usepackage{titling}
\usepackage{titlesec}
\usepackage{tocloft}
\usepackage{etoc}
\usepackage{xparse}
\renewcommand{\cfttoctitlefont}{\fontsize{24}{28}\selectfont\bfseries}\renewcommand{\cftaftertoctitle}{\hfill}\setlength{\cftbeforetoctitleskip}{0pt}\setlength{\cftaftertoctitleskip}{1.5em}

% Header/footer styling.
\usepackage{fancyhdr}
\usepackage{comment}
\newcommand{\notesauthor}{Marcus Kuo}
\newcommand{\notesinstitution}{The University of Chicago}
\newcommand{\notesyear}{2026}

\makeatletter
\renewcommand{\chaptermark}[1]{\markboth{\thechapter: #1}{}}
\makeatother

\fancypagestyle{notesbody}{%
    \fancyhf{}
	\fancyhead[L]{\itshape\nouppercase{\leftmark}}
    \fancyhead[R]{\itshape\notesauthor}
    \fancyfoot[L]{\itshape\notesyear}
    \fancyfoot[C]{\thepage}
    \renewcommand{\headrulewidth}{0pt}
    \renewcommand{\footrulewidth}{0pt}
}

\fancypagestyle{plain}{%
    \fancyhf{}
    \fancyfoot[C]{\thepage}
    \renewcommand{\headrulewidth}{0pt}
    \renewcommand{\footrulewidth}{0pt}
}

\pagestyle{notesbody}

% Section headings.
\titleformat{\section}
  {\normalfont\bfseries\large}
  {\thesection}
  {0.9em}
  {}
\titleformat{\subsection}
  {\normalfont\bfseries\large}
  {\thesubsection}
  {0.8em}
  {}
\titlespacing*{\section}{0pt}{1.4\baselineskip}{0.8\baselineskip}
\titlespacing*{\subsection}{0pt}{1.1\baselineskip}{0.55\baselineskip}

% TOC appearance.
\setcounter{tocdepth}{2}
\renewcommand{\contentsname}{Contents}

% Figure/table captions: centered, caption above, notes below by command.
\captionsetup{
    font=normalsize,
    labelfont=bf,
    textfont=normalfont,
    justification=centering,
    singlelinecheck=false,
    skip=8pt
}
\captionsetup[figure]{position=top}
\captionsetup[table]{position=top}

% Code block styling.
\usepackage{listings}
\lstdefinelanguage{Rnote}{
  language=R,
  morekeywords={mutate,summarize,filter,group_by,mean,sd,replicate,ggplot,labs,geom_density}
}
\ifnotescheatsheet
\lstdefinestyle{notescode}{
  language=Rnote,
  basicstyle=\ttfamily\small,
  keywordstyle=\color{black},
  commentstyle=\itshape\color{black},
  stringstyle=\color{black},
  showstringspaces=false,
  columns=fullflexible,
  keepspaces=true,
  aboveskip=0pt,
  belowskip=0pt,
  frame=none
}
\else
\lstdefinestyle{notescode}{
  language=Rnote,
  basicstyle=\ttfamily\small,
  keywordstyle=\color{blue!75!black},
  commentstyle=\itshape\color{teal!55!black},
  stringstyle=\color{red!70!black},
  showstringspaces=false,
  columns=fullflexible,
  keepspaces=true,
  aboveskip=0pt,
  belowskip=0pt,
  frame=none
}
\fi
\newtcblisting{rcodebox}{
  enhanced,
  breakable,
  lecrcodeboxstyle,
  listing only,
  listing options={style=notescode}
}

% Margin note box.
\makeatletter
\NewDocumentCommand{\lecnotebox}{ O{\lecboxednoteribbon} +m }{%
  \begin{tcolorbox}[
    enhanced,
    lecnoteboxstyle,
    borderline west={1.5pt}{0pt}{#1},
    width=\marginparwidth,
    before skip=0pt,
    after skip=0pt
  ]
  {%
    \begingroup
    \setstretch{1.0}%
    \raggedright
    \let\lec@orig@footnotetext\@footnotetext
    \long\def\@footnotetext##1{%
      \begingroup
      \let\footnotesize\lecboxednotefootnotefont
      \lec@orig@footnotetext{##1}%
      \endgroup
    }%
    \let\lec@orig@mpfootnotetext\@mpfootnotetext
    \long\def\@mpfootnotetext##1{%
      \begingroup
      \let\footnotesize\lecboxednotefootnotefont
      \lec@orig@mpfootnotetext{##1}%
      \endgroup
    }%
    \lecnotefont #2\par
    \endgroup
  }
  \end{tcolorbox}%
}
\NewDocumentCommand{\lecmarginparnotebox}{ +m }{%
  \begingroup
  \edef\lec@tmp{%
    \endgroup
    \noexpand\marginpar{%
      \noexpand\lecnotebox[\lecboxednoteribbon]{\unexpanded{#1}}%
    }%
  }%
  \lec@tmp
}
\NewDocumentCommand{\lecmarginnotenotebox}{ O{0pt} +m }{%
  \begingroup
  \edef\lec@tmp{%
    \endgroup
    \noexpand\marginnote{%
      \noexpand\lecnotebox[\lecboxednoteribbon]{\unexpanded{#2}}%
    }\unexpanded{[#1]}%
  }%
  \lec@tmp
}
\newbox\lec@inboxnotebox
\newdimen\lec@inboxnoteheight
\newdimen\lec@inboxnoteoffset
\newdimen\lec@inboxnoteprevbottom
\newcount\lec@inboxnotecount
\newcounter{lecboxednoteid}
\newcommand{\lecresetinboxnoteoffset}{%
  \global\lec@inboxnoteoffset=0pt\relax
  \global\lec@inboxnoteprevbottom=0pt\relax
  \global\lec@inboxnotecount=0\relax
}
\NewDocumentCommand{\lecstackedmarginnotenotebox}{ O{0pt} +m }{%
  \stepcounter{lecboxednoteid}%
  \edef\lec@inboxnoteposlabel{lecboxednote-\arabic{lecboxednoteid}}%
  \zsavepos{\lec@inboxnoteposlabel}%
  \begingroup
  \setbox\lec@inboxnotebox=\vbox{%
    \hsize=\marginparwidth
    \linewidth=\hsize
    \marginfont\raggedrightmarginnote\strut\hspace{\z@}%
    \lecnotebox[\lecboxednoteribbon]{#2}%
    \endgraf
  }%
  \global\lec@inboxnoteheight=\dimexpr\ht\lec@inboxnotebox+\dp\lec@inboxnotebox\relax
  \ifnum\zposy{\lec@inboxnoteposlabel}=0\relax
    \dimen@=\dimexpr#1+\lec@inboxnoteoffset\relax
    \global\advance\lec@inboxnoteoffset by \dimexpr\lec@inboxnoteheight+\marginparpush\relax
  \else
    \ifnum\lec@inboxnotecount=0\relax
      \dimen@=#1\relax
    \else
      \dimen@=\dimexpr\zposy{\lec@inboxnoteposlabel}sp-\lec@inboxnoteprevbottom+\marginparpush\relax
      \ifdim\dimen@<0pt\relax
        \dimen@=0pt\relax
      \fi
      \advance\dimen@ by #1\relax
    \fi
    \global\lec@inboxnoteprevbottom=\dimexpr\zposy{\lec@inboxnoteposlabel}sp-\dimen@-\lec@inboxnoteheight\relax
  \fi
  \global\advance\lec@inboxnotecount by 1\relax
  \xdef\lec@inboxnoteplacementoffset{\the\dimen@}%
  \endgroup
  \lecmarginnotenotebox[\lec@inboxnoteplacementoffset]{#2}%
}
\makeatother

% Boxed-note backends. Normal mode uses \marginpar so marginfix can prevent
% overlap and bottom-margin spillover. Notes with an explicit offset fall back
% to \marginnote as a manual escape hatch.
\NewDocumentCommand{\lecnormalboxednoteplace}{ O{0pt} +m }{%
  \ifdim#1=0pt\relax
    \lecmarginparnotebox{#2}%
  \else
    \lecmarginnotenotebox[#1]{#2}%
  \fi
}
\NewDocumentCommand{\leccheatsheetboxednoteplace}{ O{0pt} +m }{%
  \footnote{#2}%
}
\NewDocumentCommand{\lecboxednoteinboxplace}{ O{0pt} +m }{%
  \ifnotescheatsheet
    \leccheatsheetboxednoteplace[#1]{#2}%
  \else
    \lecstackedmarginnotenotebox[#1]{#2}%
  \fi
}
\let\lecboxednoteplace\lecnormalboxednoteplace
\ifnotescheatsheet
\let\lecboxednoteplace\leccheatsheetboxednoteplace
\fi

% Notes triggered in math mode or other inner boxes are queued until we hit
% a stable text boundary where the chosen backend can place them safely.
\makeatletter
\newcommand{\lecdeferrednotes}{}
\NewDocumentCommand{\lecdefernote}{ O{0pt} +m }{%
  \g@addto@macro\lecdeferrednotes{\lecboxednoteplace[#1]{#2}}%
}
\newcommand{\lecflushdeferrednotes}{%
  \lecdeferrednotes
  \global\let\lecdeferrednotes\@empty
}
\AtEndDocument{\lecflushdeferrednotes}
\pretocmd{\appendix}{\lecflushdeferrednotes}{}{}
\makeatother

\NewDocumentCommand{\boxednote}{ O{0pt} +m }{%
  \ifmmode
    \lecdefernote[#1]{#2}%
  \else\iflecinbox
    \lecboxednoteinboxplace[#1]{#2}%
  \else\ifinner
    \lecdefernote[#1]{#2}%
  \else
    \lecflushdeferrednotes
    \lecboxednoteplace[#1]{#2}%
  \fi\fi\fi
}

% Margin figure helper: caption (with number) above, image nearly full width.
% Usage: \marginimage[Caption][fig:label][Figure note.]{filename}
% - `filename` is relative to `images/`
% - Omit caption/label/note by leaving the optional args empty
\NewDocumentCommand{\marginimage}{ O{} o O{} m }{%
  \boxednote{%
    \begingroup
    \centering
    \if\relax\detokenize{#1}\relax
      \IfNoValueF{#2}{%
        \if\relax\detokenize{#2}\relax\else
          \refstepcounter{figure}%
          \label{#2}%
        \fi
      }%
    \else
      \refstepcounter{figure}%
      \IfNoValueF{#2}{%
        \if\relax\detokenize{#2}\relax\else\label{#2}\fi
      }%
      {\footnotesize\bfseries \figurename~\thefigure.\ \itshape #1\par}%
      \vspace{4pt}%
    \fi
    \includegraphics[width=\linewidth]{images/#4}%
    \if\relax\detokenize{#3}\relax\else
      \fignote{#3}%
    \fi
    \endgroup
  }%
}
\ifnotescheatsheet
\RenewDocumentCommand{\marginimage}{ O{} o O{} m }{}%
\fi

% Notes under tables/figures.
\newcommand{\fignote}[1]{\par\vspace{0.25em}{\footnotesize\textit{Notes:} #1\par}}
\newcommand{\tabnote}[1]{\par\vspace{0.25em}{\footnotesize\textit{Notes:} #1\par}}

% Line-break penalties: prevent mid-expression line breaks.
\relpenalty=10000
\binoppenalty=10000

% =====================================================================
% MATH SHORTHAND MACROS
% =====================================================================
\newcommand*\N{\mathbb{N}}
\newcommand*\R{\mathbb{R}}
\newcommand*\Q{\mathbb{Q}}
\newcommand*\Z{\mathbb{Z}}
\newcommand*\E{\mathbb{E}}
\let\oldd\d
\renewcommand*\d{\mathop{}\!\mathrm{d}}
\let\oldL\L
\let\oldP\P
\renewcommand*\P{\mathbb{P}}
\renewcommand*\L{\mathbb{L}}
\let\O\varnothing
\let\oldemptyset\emptyset
\let\emptyset\varnothing
\newcommand{\ind}{\perp\!\!\!\perp}

\DeclareMathOperator*{\argmax}{arg\,max}
\DeclareMathOperator*{\argmin}{arg\,min}
\DeclareMathOperator*{\plim}{plim}
\DeclareMathOperator*{\esssup}{ess\,sup}
\DeclareMathOperator*{\essinf}{ess\,inf}
\DeclareMathOperator{\opvec}{vec}
\DeclareMathOperator{\Var}{Var}
\DeclareMathOperator{\Cov}{Cov}
\DeclareMathOperator{\Corr}{Corr}
\DeclareMathOperator{\Bias}{Bias}
\DeclareMathOperator{\MSE}{MSE}

% =====================================================================
% CROSS-REFERENCE MACROS
% =====================================================================
\newcommand{\thmref}[1]{Theorem~\ref{thm:#1}}
\newcommand{\lemref}[1]{Lemma~\ref{thm:#1}}
\newcommand{\propref}[1]{Proposition~\ref{thm:#1}}
\newcommand{\defref}[1]{Definition~\ref{thm:#1}}
\newcommand{\corref}[1]{Corollary~\ref{thm:#1}}
\newcommand{\exref}[1]{Example~\ref{thm:#1}}

% =====================================================================
% THEOREM ENVIRONMENTS
% =====================================================================

% Shared counter: chapter.section.number (e.g., 1.2.3)
\newcounter{thmcounter}[section]
\renewcommand{\thethmcounter}{\thesection.\arabic{thmcounter}}

% Colors for theorem boxes.
% User palette (hex):
%  - theorem/proposition: #ccd5ae
%  - lemma/corollary:     #ecf1d3
%  - assumption:          #ddb892
%  - definition/notation: #fefae0
%  - example:             #ffd9c7
\definecolor{lecthm}{HTML}{CCD5AE}
\definecolor{leclem}{HTML}{ECF1D3}
\definecolor{lecassump}{HTML}{DDB892}
\definecolor{lecdef}{HTML}{FEFAE0}
\definecolor{lecex}{HTML}{FFD9C7}

% Main box fill shade: lower = lighter (more white mixed in),
% higher = darker (closer to the base color).
\newcommand{\lecmainmix}{32}

% Title bars are mixed with neutral gray: lower = less saturated.
\newcommand{\lectitlebarmix}{92}
\newcommand{\lecframemix}{95}
\ifnotescheatsheet
\colorlet{lecthm}{white}
\colorlet{leclem}{white}
\colorlet{lecassump}{white}
\colorlet{lecdef}{white}
\colorlet{lecex}{white}
\renewcommand{\lecmainmix}{100}
\renewcommand{\lectitlebarmix}{100}
\renewcommand{\lecframemix}{0}
\fi

% Proof-end symbol mechanism.
\newcommand{\currentproofmark}{\ensuremath{\square}}
\newcommand{\setproofmark}[1]{\gdef\currentproofmark{#1}}
\newcommand{\resetproofmark}{\setproofmark{\ensuremath{\square}}}
\DeclareRobustCommand{\leciconqed}[1]{%
  \ifmmode
    \text{\raisebox{-0.1ex}{#1}}%
  \else
    \raisebox{-0.1ex}{#1}%
  \fi
}

% Default proof label (auto-updated by theorem-like environments).
\newcommand{\lecproofname}{Proof}
\newcommand{\setlecproofname}[1]{%
  \begingroup
  \edef\lecprooftmp{#1}%
  \global\let\lecproofname\lecprooftmp
  \endgroup
}
\newcommand{\resetlecproofname}{\setlecproofname{Proof}}
\newcommand{\theoremproofstyle}{\setproofmark{\leciconqed{\faKiwiBird}}}
\newcommand{\propositionproofstyle}{\setproofmark{\leciconqed{\faKiwiBird}}}
\newcommand{\lemmaproofstyle}{\setproofmark{\leciconqed{\faLockOpen}}}
\newcommand{\corollaryproofstyle}{\setproofmark{\leciconqed{\faKiwiBird}}}
\newcommand{\exampleproofstyle}{\setproofmark{\leciconqed{\faFrog}}}
\newcommand{\remarkproofstyle}{\setproofmark{\leciconqed{\faReadme}}}

% Base tcolorbox style for theorem environments.
\tcbset{
  thmbase/.style={
    enhanced,
    breakable,
    boxrule=0.32pt,
    arc=0.6mm,
    left=6pt,
    right=6pt,
    top=4pt,
    bottom=4pt,
    before skip=4pt,
    after skip=5pt,
    fonttitle=\bfseries\fontsize{9.75}{11.7}\selectfont,
    coltitle=black,
    titlerule=0pt,
    toptitle=2pt,
    bottomtitle=2pt,
    separator sign none,
    description delimiters parenthesis,
    terminator sign={.},
  },
	  thmbluestyle/.style={
	    thmbase,
	    colback=lecthm!8!white,
	    colframe=lecthm!86!black,
	    colbacktitle=lecthm!90!gray,
	  },
	  thmbgstyle/.style={
	    thmbase,
	    colback=leclem!8!white,
	    colframe=leclem!82!black,
	    colbacktitle=leclem!90!gray,
	  },
	  thmyellowstyle/.style={
	    thmbase,
	    colback=lecdef!7!white,
	    colframe=lecdef!85!black,
	    colbacktitle=lecdef!92!gray,
	  },
	  thmpinkstyle/.style={
	    thmbase,
	    colback=lecex!7!white,
	    colframe=lecex!80!black,
	    colbacktitle=lecex!88!gray,
	  },
	  thmassumpstyle/.style={
	    thmbase,
	    colback=lecassump!7!white,
	    colframe=lecassump!82!black,
	    colbacktitle=lecassump!88!gray,
	  },
	}

% Boxed theorem environments via newtcbtheorem.
% Usage: \begin{theorem}{Name}{label} ... \end{theorem}
%   - Name appears in parentheses after "Theorem 1.2.3"
%   - Empty name: \begin{theorem}{}{label}
%   - Label is prefixed with "thm:" automatically

\newtcbtheorem[use counter=thmcounter, number freestyle={\noexpand\thesection.\noexpand\arabic{thmcounter}}]{lectheorem}{Theorem}{
  thmbluestyle,
  before upper={
    \lecsetboxednoteribbon{lecthm!90!gray}
    % theorem, proposition, corollary: faKiwiBird # the kiwibird emoji is cute
    \theoremproofstyle
    \setlecproofname{Proof of Theorem~\thethmcounter}
  },
}{thm}

\newtcbtheorem[use counter=thmcounter, number freestyle={\noexpand\thesection.\noexpand\arabic{thmcounter}}]{lecproposition}{Proposition}{
  thmbluestyle,
  before upper={
    \lecsetboxednoteribbon{lecthm!90!gray}
    % theorem, proposition, corollary: faKiwiBird # the kiwibird emoji is cute
    \propositionproofstyle
    \setlecproofname{Proof of Proposition~\thethmcounter}
  },
}{thm}

\newtcbtheorem[use counter=thmcounter, number freestyle={\noexpand\thesection.\noexpand\arabic{thmcounter}}]{leclemma}{Lemma}{
  thmbgstyle,
  before upper={
    \lecsetboxednoteribbon{leclem!90!gray}
    % lemma: faLockOpen # lemma and lightbulb start with L, and lemma proofs unlock tools for later use.
    \lemmaproofstyle
    \setlecproofname{Proof of Lemma~\thethmcounter}
  },
}{thm}

\newtcbtheorem[use counter=thmcounter, number freestyle={\noexpand\thesection.\noexpand\arabic{thmcounter}}]{leccorollary}{Corollary}{
  thmbgstyle,
  before upper={
    \lecsetboxednoteribbon{leclem!90!gray}
    % theorem, proposition, corollary: faKiwiBird # the kiwibird emoji is cute
    \corollaryproofstyle
    \setlecproofname{Proof of Corollary~\thethmcounter}
  },
}{thm}

\newtcbtheorem[use counter=thmcounter, number freestyle={\noexpand\thesection.\noexpand\arabic{thmcounter}}]{lecdefinition}{Definition}{
  thmyellowstyle,
  before upper={
    \lecsetboxednoteribbon{lecdef!92!gray}
  },
}{thm}

\newtcbtheorem[use counter=thmcounter, number freestyle={\noexpand\thesection.\noexpand\arabic{thmcounter}}]{lecexample}{Example}{
  thmpinkstyle,
  before upper={
    \lecsetboxednoteribbon{lecex!88!gray}
    % example: faFrog
    \exampleproofstyle
  },
  overlay unbroken and last={
    \node[anchor=south east, inner sep=6pt] at (frame.south east) {\leciconqed{\faFrog}};
  },
}{thm}

\newtcbtheorem[use counter=thmcounter, number freestyle={\noexpand\thesection.\noexpand\arabic{thmcounter}}]{lecassumption}{Assumption}{
  thmassumpstyle,
  before upper={
    \lecsetboxednoteribbon{lecassump!88!gray}
  },
}{thm}

\newtcbtheorem[use counter=thmcounter, number freestyle={\noexpand\thesection.\noexpand\arabic{thmcounter}}]{lecnotation}{Notation}{
  thmyellowstyle,
  before upper={
    \lecsetboxednoteribbon{lecdef!92!gray}
  },
}{thm}

% Wrapper environments: optional title, optional label.
%   - \begin{theorem}[Title] ... \end{theorem} (auto-label; use \label{...} only if you want to refer)
%   - \begin{theorem}[Title][mylabel] ... \end{theorem} (sets \label{thm:mylabel})
\newcounter{lecautolabel}
\newif\iflecinbox
\makeatletter
\newcommand{\lecsetthmcurrentlabel}{%
  \protected@edef\@currentlabel{\thethmcounter}%
}
\makeatother
\NewDocumentEnvironment{theorem}{ O{} o }{%
  \global\lecinboxtrue
  \lecresetinboxnoteoffset
  \IfNoValueTF{#2}{%
    \stepcounter{lecautolabel}%
    \begin{lectheorem}{#1}{auto-\arabic{lecautolabel}}%
    \lecsetthmcurrentlabel
  }{%
    \begin{lectheorem}{#1}{#2}%
    \lecsetthmcurrentlabel
  }%
}{\end{lectheorem}\global\lecinboxfalse\lecflushdeferrednotes\lecresetboxednoteribbon}
\NewDocumentEnvironment{proposition}{ O{} o }{%
  \global\lecinboxtrue
  \lecresetinboxnoteoffset
  \IfNoValueTF{#2}{%
    \stepcounter{lecautolabel}%
    \begin{lecproposition}{#1}{auto-\arabic{lecautolabel}}%
    \lecsetthmcurrentlabel
  }{%
    \begin{lecproposition}{#1}{#2}%
    \lecsetthmcurrentlabel
  }%
}{\end{lecproposition}\global\lecinboxfalse\lecflushdeferrednotes\lecresetboxednoteribbon}
\NewDocumentEnvironment{lemma}{ O{} o }{%
  \global\lecinboxtrue
  \lecresetinboxnoteoffset
  \IfNoValueTF{#2}{%
    \stepcounter{lecautolabel}%
    \begin{leclemma}{#1}{auto-\arabic{lecautolabel}}%
    \lecsetthmcurrentlabel
  }{%
    \begin{leclemma}{#1}{#2}%
    \lecsetthmcurrentlabel
  }%
}{\end{leclemma}\global\lecinboxfalse\lecflushdeferrednotes\lecresetboxednoteribbon}
\NewDocumentEnvironment{corollary}{ O{} o }{%
  \global\lecinboxtrue
  \lecresetinboxnoteoffset
  \IfNoValueTF{#2}{%
    \stepcounter{lecautolabel}%
    \begin{leccorollary}{#1}{auto-\arabic{lecautolabel}}%
    \lecsetthmcurrentlabel
  }{%
    \begin{leccorollary}{#1}{#2}%
    \lecsetthmcurrentlabel
  }%
}{\end{leccorollary}\global\lecinboxfalse\lecflushdeferrednotes\lecresetboxednoteribbon}
\NewDocumentEnvironment{definition}{ O{} o }{%
  \global\lecinboxtrue
  \lecresetinboxnoteoffset
  \IfNoValueTF{#2}{%
    \stepcounter{lecautolabel}%
    \begin{lecdefinition}{#1}{auto-\arabic{lecautolabel}}%
    \lecsetthmcurrentlabel
  }{%
    \begin{lecdefinition}{#1}{#2}%
    \lecsetthmcurrentlabel
  }%
}{\end{lecdefinition}\global\lecinboxfalse\lecflushdeferrednotes\lecresetboxednoteribbon}
\NewDocumentEnvironment{example}{ O{} o }{%
  \global\lecinboxtrue
  \lecresetinboxnoteoffset
  \IfNoValueTF{#2}{%
    \stepcounter{lecautolabel}%
    \begin{lecexample}{#1}{auto-\arabic{lecautolabel}}%
    \lecsetthmcurrentlabel
  }{%
    \begin{lecexample}{#1}{#2}%
    \lecsetthmcurrentlabel
  }%
}{\end{lecexample}\global\lecinboxfalse\lecflushdeferrednotes\lecresetboxednoteribbon}
\NewDocumentEnvironment{assumption}{ O{} o }{%
  \global\lecinboxtrue
  \lecresetinboxnoteoffset
  \IfNoValueTF{#2}{%
    \stepcounter{lecautolabel}%
    \begin{lecassumption}{#1}{auto-\arabic{lecautolabel}}%
    \lecsetthmcurrentlabel
  }{%
    \begin{lecassumption}{#1}{#2}%
    \lecsetthmcurrentlabel
  }%
}{\end{lecassumption}\global\lecinboxfalse\lecflushdeferrednotes\lecresetboxednoteribbon}
\NewDocumentEnvironment{notation}{ O{} o }{%
  \global\lecinboxtrue
  \lecresetinboxnoteoffset
  \IfNoValueTF{#2}{%
    \stepcounter{lecautolabel}%
    \begin{lecnotation}{#1}{auto-\arabic{lecautolabel}}%
    \lecsetthmcurrentlabel
  }{%
    \begin{lecnotation}{#1}{#2}%
    \lecsetthmcurrentlabel
  }%
}{\end{lecnotation}\global\lecinboxfalse\lecflushdeferrednotes\lecresetboxednoteribbon}

% Unboxed environments: remark and aside (amsthm, sharing thmcounter).
\theoremstyle{remark}
\newtheorem{remark}[thmcounter]{Remark}
\newtheorem{aside}[thmcounter]{Aside}
\newtheorem{conj}[thmcounter]{Conjecture}
\let\lecoriginalremark\remark
\let\endlecoriginalremark\endremark
\RenewDocumentEnvironment{remark}{ o }{%
  \remarkproofstyle
  \pushQED{\qed}%
  \IfNoValueTF{#1}{\lecoriginalremark}{\lecoriginalremark[#1]}%
}{%
  \popQED
  \endlecoriginalremark
  \resetproofmark
}

% Override proof environment to use \currentproofmark as QED symbol.
\makeatletter
\renewenvironment{proof}[1][\lecproofname]{%
  \par\pushQED{\qed}%
  \normalfont \topsep6\p@\@plus6\p@\relax
  \trivlist
  \item[\hskip\labelsep\itshape #1\@addpunct{.}]\ignorespaces
}{%
  \popQED\endtrivlist\@endpefalse
  \resetlecproofname
  \resetproofmark
}
\makeatother
\renewcommand{\qedsymbol}{\currentproofmark}

% =====================================================================
% TITLE PAGE MACRO
% =====================================================================
\newcommand{\makelectitle}[3]{%
  \begin{titlepage}
    \thispagestyle{empty}
    \begin{tikzpicture}[remember picture,overlay]
      % Title block (positions tuned to match the reference layout).
      \node[
        anchor=north,
        inner sep=0pt,
        % Use some of the outer margin for long titles, but keep within the page.
        text width=\dimexpr\textwidth+\marginparsep+\marginparwidth-0.60in\relax,
        align=center
      ] at ([yshift=-3.55in]current page.north) {%
        \setstretch{1.0}%
        \fontsize{24}{28}\selectfont\bfseries
        \spaceskip=0.35em plus 0.04em minus 0.015em\relax
        \xspaceskip=\spaceskip
        \hyphenchar\font=-1\relax
        \centering
        % Prefer non-hyphenated breaks on long titles.
        \pretolerance=10000\relax
        \tolerance=10000\relax
        \emergencystretch=2em\relax
        \hyphenpenalty=10000\relax
        \exhyphenpenalty=10000\relax
        #1%
      };

      \node[
        anchor=north,
        inner sep=0pt,
        text width=\textwidth,
        align=center
      ] at ([yshift=-5.05in]current page.north) {%
        \fontsize{17}{20}\selectfont\itshape #2%
      };

      \node[
        anchor=north,
        inner sep=0pt,
        text width=\textwidth,
        align=center
      ] at ([yshift=-6.40in]current page.north) {%
        \fontsize{14}{16}\selectfont\scshape #3%
      };
    \end{tikzpicture}
  \end{titlepage}
  \clearpage
}

% =====================================================================
% CHAPTER OPENER MACRO
% =====================================================================
\newlength{\chaptextwid}
\newcommand{\chapteropening}[1]{%
  \thispagestyle{plain}
  \vspace*{0.10in}
  \noindent
  \setlength{\chaptextwid}{\textwidth}%
  \begin{tikzpicture}[x=1in,y=1in]
    \pgfmathsetmacro{\rr}{\chaptextwid/1in}
    \pgfmathsetmacro{\chapruletop}{0.10}
    \pgfmathsetmacro{\chaprulebottom}{-0.50}
    \pgfmathsetmacro{\chaprulethickness}{0.015}
    \pgfmathsetmacro{\chapcontenty}{(\chapruletop+\chaprulebottom)/2}
    \path[fill=black] (0,\chapruletop) rectangle (\rr,\chapruletop+\chaprulethickness);
    \node[
      anchor=west,
      font=\fontsize{19}{21}\selectfont\bfseries,
      inner sep=0pt,
      text width=\dimexpr\chaptextwid-0.62in\relax,
      align=left
    ] at (0,\chapcontenty) {#1};
    \node[anchor=east,font=\fontsize{30}{30}\selectfont\bfseries] at (\rr,\chapcontenty) {\thechapter};
    \path[fill=black] (0,\chaprulebottom-\chaprulethickness) rectangle (\rr,\chaprulebottom);
  \end{tikzpicture}
  \vspace{0.26in}
  {\parindent=0pt
   \etocsettocstyle{}{}
   \etocsetnexttocdepth{section}
   \etocsetstyle{section}
     {}
     {\makebox[2.2em][l]{\etocnumber}\etocname\nobreak\leaders\hbox to 0.7em{.\hss}\hfill\etocpage\par}
     {}
     {}
   \localtableofcontents
  }
  \vspace{0.35in}
}

% Convenience wrapper: manual chapter step + styled opener.
\newcommand{\styledchapter}[1]{%
  \lecflushdeferrednotes
  \cleardoublepage
  \refstepcounter{chapter}%
  \setcounter{section}{0}%
  \setcounter{subsection}{0}%
  \setcounter{equation}{0}%
  \markboth{\thechapter: #1}{}%
  \addcontentsline{toc}{chapter}{\protect\numberline{\thechapter}#1}%
  \chapteropening{#1}%
}
